\section{课程整体结构}

本讲最后给出课程的完整地图。课程以“效率”为主线，把内容分成五个大块：

\begin{table}[H]
\centering
\begin{tabular}{p{2.2cm}p{4.1cm}p{6.8cm}}
\toprule
\textbf{模块} & \textbf{关心的问题} & \textbf{你会学到什么} \\
\midrule
Basics & 从零搭出一个基本 pipeline & tokenization、Transformer、训练循环 \\
Systems & 如何把硬件压榨到极致 & kernels、并行训练、推理优化 \\
Scaling laws & 如何用小实验预测大规模表现 & compute-optimal、超参外推、预算分配 \\
Data & 什么数据值得训练、如何处理数据 & 采集、清洗、过滤、去重、评估 \\
Alignment & 如何让 base model 真正可用 & SFT、DPO、GRPO、风格和安全 \\
\bottomrule
\end{tabular}
\end{table}

这门课的核心思想是：你永远不是在抽象真空里做研究，而是在资源约束下做选择。
因此，课程的每一部分都会反复回到“效率”这个关键词。

